\begin{thebibliography}{99}
\bibitem{bib-bib1} 
\par\noindent    Rajeev Alur, Caleb Stanford, and Christopher Watson.  A robust theory of series parallel graphs.  Proc. ACM Program. Lang., 7(POPL):1058–1088, 2023.  doi:10.1145/3571230.  

\bibitem{bib-bib2} 
\par\noindent    Amazigh Amrane, Hugo Bazille, Emily Clement, and Uli Fahrenberg.  Languages of higher-dimensional timed automata.  In Lars Michael Kristensen and Jan Martijn E. M. van der Werf, editors, Application and Theory of Petri Nets and Concurrency - 45th International Conference, PETRI NETS 2024, Geneva, Switzerland, June 26-28, 2024, Proceedings, volume 14628 of Lecture Notes in Computer Science, pages 197–219. Springer, 2024.  doi:10.1007/978-3-031-61433-0\_10.  

\bibitem{bib-bib3} 
\par\noindent    Amazigh Amrane, Hugo Bazille, Emily Clement, Uli Fahrenberg, and Philipp Schlehuber-Caissier.  Higher-dimensional timed automata for real-time concurrency.  In Andreas Rauh, Bernd Finkbeiner, and Paul Kröger, editors, Design and Verification of Cyber-Physical Systems. Springer, 2025.  To be published. https://arxiv.org/abs/2401.17444.  

\bibitem{bib-bib4} 
\par\noindent    Amazigh Amrane, Hugo Bazille, Emily Clement, Uli Fahrenberg, and Krzysztof Ziemianski.  Presenting interval pomsets with interfaces.  In Uli Fahrenberg, Wesley Fussner, and Roland Glück, editors, Relational and Algebraic Methods in Computer Science - 21st International Conference, RAMiCS 2024, Prague, Czech Republic, August 19-22, 2024, Proceedings, volume 14787 of Lecture Notes in Computer Science, pages 28–45. Springer, 2024.  doi:10.1007/978-3-031-68279-7\_3.  

\bibitem{bib-bib5} 
\par\noindent    Amazigh Amrane, Hugo Bazille, Uli Fahrenberg, and Marie Fortin.  Logic and languages of higher-dimensional automata.  In International Conference on Developments in Language Theory, pages 51–67. Springer, 2024.  doi:10.1007/978-3-031-66159-4\_5.  

\bibitem{bib-bib6} 
\par\noindent    Amazigh Amrane, Hugo Bazille, Uli Fahrenberg, Loïc Hélouët, and Philipp Schlehuber-Caissier.  Petri nets and higher-dimensional automata, 2025.  Accepted for PETRI NETS.  doi:10.48550/arXiv.2502.02354.  

\bibitem{bib-bib7} 
\par\noindent    Amazigh Amrane, Hugo Bazille, Uli Fahrenberg, and Krzysztof Ziemiański.  Closure and decision properties for higher-dimensional automata.  In Erika Ábrahám, Clemens Dubslaff, and Silvia Lizeth Tapia Tarifa, editors, ICTAC, volume 14446 of Lecture Notes in Computer Science, pages 295–312. Springer, 2023.  doi:10.1007/978-3-031-47963-2\_18.  

\bibitem{bib-bib8} 
\par\noindent    Amazigh Amrane, Hugo Bazille, Uli Fahrenberg, and Krzysztof Ziemiański.  Closure and decision properties for higher-dimensional automata.  Theoretical Computer Science, 1036:115156, 2025.  doi:10.1016/j.tcs.2025.115156.  

\bibitem{bib-bib9} 
\par\noindent    Marek A. Bednarczyk.  Categories of Asynchronous Systems.  PhD thesis, University of Sussex, UK, 1987.  

\bibitem{bib-bib10} 
\par\noindent    J Richard Büchi.  Symposium on decision problems: On a decision method in restricted second order arithmetic.  In Studies in Logic and the Foundations of Mathematics, volume 44, pages 1–11. Elsevier, 1966.  

\bibitem{bib-bib11} 
\par\noindent    Olivier Carton, Dominique Perrin, and Jean-Eric Pin.  Automata and semigroups recognizing infinite words.  Logic and automata, 2:133–168, 2008.  

\bibitem{bib-bib12} 
\par\noindent    Volker Diekert and Anca Muscholl.  Deterministic asynchronous automata for infinite traces.  In Annual Symposium on Theoretical Aspects of Computer Science, pages 617–628. Springer, 1993.  doi:10.1007/3-540-56503-5\_61.  

\bibitem{bib-bib13} 
\par\noindent    Jérémy Dubut.  Trees in partial higher dimensional automata.  In Mikołaj Bojańczyk and Alex Simpson, editors, FOSSACS, volume 11425 of Lecture Notes in Computer Science, pages 224–241. Springer, 2019.  doi:10.1007/978-3-030-17127-8\_13.  

\bibitem{bib-bib14} 
\par\noindent    Alexandre Duret-Lutz, Etienne Renault, Maximilien Colange, Florian Renkin, Alexandre Gbaguidi Aisse, Philipp Schlehuber-Caissier, Thomas Medioni, Antoine Martin, Jérôme Dubois, Clément Gillard, and Henrich Lauko.  From Spot 2.0 to Spot 2.10: What’s new?  In 34th International Conference on Computer Aided Verification (CAV’22), volume 13372 of Lecture Notes in Computer Science, pages 174–187. Springer, 2022.  doi:10.1007/978-3-031-13188-2\_9.  

\bibitem{bib-bib15} 
\par\noindent    Werner Ebinger and Anca Muscholl.  Logical definability on infinite traces.  Theoretical Computer Science, 154(1):67–84, 1996.  doi:10.1016/0304-3975(95)00130-1.  

\bibitem{bib-bib16} 
\par\noindent    Uli Fahrenberg.  Higher-dimensional timed automata.  In Alessandro Abate, Antoine Girard, and Maurice Heemels, editors, ADHS, volume 51 of IFAC-PapersOnLine, pages 109–114. Elsevier, 2018.  doi:10.1016/j.ifacol.2018.08.019.  

\bibitem{bib-bib17} 
\par\noindent    Uli Fahrenberg.  Higher-dimensional timed and hybrid automata.  Leibniz Transactions on Embedded Systems, 8(2):03:1–03:16, 2022.  doi:10.4230/LITES.8.2.3.  

\bibitem{bib-bib18} 
\par\noindent    Uli Fahrenberg, Christian Johansen, Georg Struth, and Krzysztof Ziemiański.  Languages of higher-dimensional automata.  Mathematical Structures in Computer Science, 31(5):575–613, 2021.  doi:10.1017/S0960129521000293.  

\bibitem{bib-bib19} 
\par\noindent    Uli Fahrenberg, Christian Johansen, Georg Struth, and Krzysztof Ziemiański.  A Kleene theorem for higher-dimensional automata.  In Bartek Klin, Sławomir Lasota, and Anca Muscholl, editors, CONCUR, volume 243 of Leibniz International Proceedings in Informatics, pages 29:1–29:18. Schloss Dagstuhl – Leibniz-Zentrum für Informatik, 2022.  doi:10.4230/LIPIcs.CONCUR.2022.29.  

\bibitem{bib-bib20} 
\par\noindent    Uli Fahrenberg, Christian Johansen, Georg Struth, and Krzysztof Ziemiański.  Posets with interfaces as a model for concurrency.  Information and Computation, 285(B):104914, 2022.  doi:10.1016/j.ic.2022.104914.  

\bibitem{bib-bib21} 
\par\noindent    Uli Fahrenberg and Axel Legay.  Partial higher-dimensional automata.  In Lawrence S. Moss and Pawel Sobocinski, editors, CALCO, volume 35 of LIPIcs, pages 101–115. Schloss Dagstuhl – Leibniz-Zentrum für Informatik, 2015.  doi:10.4230/LIPIcs.CALCO.2015.101.  

\bibitem{bib-bib22} 
\par\noindent    Uli Fahrenberg and Krzysztof Ziemiański.  A Myhill-Nerode theorem for higher-dimensional automata.  In Luís Gomes and Robert Lorenz, editors, PETRI NETS, volume 13929 of Lecture Notes in Computer Science, pages 167–188. Springer, 2023.  doi:10.1007/978-3-031-33620-1\_9.  

\bibitem{bib-bib23} 
\par\noindent    Jean Fanchon and Rémi Morin.  Pomset languages of finite step transition systems.  In Giuliana Franceschinis and Karsten Wolf, editors, PETRI NETS, volume 5606 of Lecture Notes in Computer Science, pages 83–102. Springer, 2009.  doi:10.1007/978-3-642-02424-5\_7.  

\bibitem{bib-bib24} 
\par\noindent    Berndt Farwer.  omega-automata.  In Erich Grädel, Wolfgang Thomas, and Thomas Wilke, editors, Automata, Logics, and Infinite Games: A Guide to Current Research [outcome of a Dagstuhl seminar, February 2001], volume 2500 of Lecture Notes in Computer Science, pages 3–20. Springer, 2001.  doi:10.1007/3-540-36387-4\_1.  

\bibitem{bib-bib25} 
\par\noindent    Peter C. Fishburn.  Intransitive indifference with unequal indifference intervals.  Journal of Mathematical Psychology, 7(1):144–149, 1970.  doi:10.1016/0022-2496(70)90062-3.  

\bibitem{bib-bib26} 
\par\noindent    Paul Gastin, Antoine Petit, and Wiesław Zielonka.  An extension of Kleene’s and Ochmański’s theorems to infinite traces.  Theoretical Computer Science, 125(2):167–204, 1994.  doi:10.1016/0304-3975(94)90254-2.  

\bibitem{bib-bib27} 
\par\noindent    Blaise Genest, Dietrich Kuske, and Anca Muscholl.  A Kleene theorem and model checking algorithms for existentially bounded communicating automata.  Information and Computation, 204(6):920–956, 2006.  doi:10.1016/j.ic.2006.01.005.  

\bibitem{bib-bib28} 
\par\noindent    Jay L Gischer.  The equational theory of pomsets.  Theoretical Computer Science, 61(2-3):199–224, 1988.  doi:10.1016/0304-3975(88)90124-7.  

\bibitem{bib-bib29} 
\par\noindent    Jan Grabowski.  On partial languages.  Fundamentae Informatica, 4(2):427, 1981.  

\bibitem{bib-bib30} 
\par\noindent    Kenneth Kunen.  Set Theory: An Introduction to Independence Proofs.  North-Holland, 1980.  

\bibitem{bib-bib31} 
\par\noindent    Dietrich Kuske.  Infinite series-parallel posets: Logic and languages.  In Ugo Montanari, José D. P. Rolim, and Emo Welzl, editors, Automata, Languages and Programming, 27th International Colloquium, ICALP 2000, Geneva, Switzerland, July 9-15, 2000, Proceedings, volume 1853 of Lecture Notes in Computer Science, pages 648–662. Springer, 2000.  doi:10.1007/3-540-45022-X\_55.  

\bibitem{bib-bib32} 
\par\noindent    K. Lodaya and P. Weil.  A Kleene iteration for parallelism.  In Foundations of Software Technology and Theoretical Computer Science, pages 355–366, 1998.  

\bibitem{bib-bib33} 
\par\noindent    K. Lodaya and P. Weil.  Series-parallel posets: algebra, automata and languages.  In STACS’98, volume 1373 of Lecture Notes in Computer Science, pages 555–565. Springer, 1998.  doi:10.1007/BFB0028590.  

\bibitem{bib-bib34} 
\par\noindent    K. Lodaya and P. Weil.  Series-parallel languages and the bounded-width property.  Theoretical Computer Science, 237(1–2):347–380, 2000.  doi:10.1016/S0304-3975(00)00031-1.  

\bibitem{bib-bib35} 
\par\noindent    K. Lodaya and P. Weil.  Rationality in algebras with a series operation.  Information and Computation, 171:269–293, 2001.  doi:10.1006/INCO.2001.3077.  

\bibitem{bib-bib36} 
\par\noindent    Robert McNaughton.  Testing and generating infinite sequences by a finite automaton.  Information and control, 9(5):521–530, 1966.  doi:10.1016/S0019-9958(66)80013-X.  

\bibitem{bib-bib37} 
\par\noindent    David E Muller.  Infinite sequences and finite machines.  In Proceedings of the Fourth Annual Symposium on Switching Circuit Theory and Logical Design (swct 1963), pages 3–16. IEEE Computer Society, 1963.  doi:10.1109/SWCT.1963.8.  

\bibitem{bib-bib38} 
\par\noindent    Luc Passemard, Amazigh Amrane, and Uli Fahrenberg.  Higher-dimensional automata : Extension to infinite tracks.  CoRR, abs/2503.07881, 2025.  doi:10.48550/arXiv.2503.07881.  

\bibitem{bib-bib39} 
\par\noindent    Dominique Perrin and Jean-Eric Pin.  Infinite words - automata, semigroups, logic and games, volume 141 of Pure and applied mathematics series.  Elsevier Morgan Kaufmann, 2004.  

\bibitem{bib-bib40} 
\par\noindent    Vaughan R. Pratt.  Modeling concurrency with partial orders.  Journal of Parallel Programming, 15(1):33–71, 1986.  doi:10.1007/BF01379149.  

\bibitem{bib-bib41} 
\par\noindent    Vaughan R. Pratt.  Modeling concurrency with geometry.  In POPL, pages 311–322, New York City, 1991. ACM Press.  doi:10.1145/99583.99625.  

\bibitem{bib-bib42} 
\par\noindent    Rob J. van Glabbeek.  Bisimulations for higher dimensional automata.  Email message, June 1991.  URL: http://theory.stanford.edu/˜rvg/hda.  

\bibitem{bib-bib43} 
\par\noindent    Rob J. van Glabbeek.  On the expressiveness of higher dimensional automata.  Theoretical Computer Science, 356(3):265–290, 2006.  doi:10.1016/j.tcs.2006.02.012.  

\bibitem{bib-bib44} 
\par\noindent    Glynn Winskel.  Event structures.  In W. Brauer, W. Reisig, and G. Rozenberg, editors, Petri Nets: Applications and Relationships to Other Models of Concurrency, pages 325–392, Berlin, Heidelberg, 1987. Springer Berlin Heidelberg.  

\bibitem{bib-bib45} 
\par\noindent    Glynn Winskel and Mogens Nielsen.  Models for concurrency.  In Samson Abramsky, Dov M. Gabbay, and Thomas S.E. Maibaum, editors, Handbook of Logic in Computer Science, volume 4. Clarendon Press, Oxford, 1995.  

\bibitem{bib-bib46} 
\par\noindent    Wiesław Zielonka.  Notes on finite asynchronous automata.  RAIRO – Informatique Théorique et Applications, 21(2):99–135, 1987.  doi:10.1051/ita/1987210200991.  

\end{thebibliography}
